\chapter{\IfLanguageName{dutch}{Technische uitvoering}{Technical implementation}}
\label{ch:technische-uitvoering}

De voorgaande fases resulteerden in concrete shortlists van de drie variabele componenten die samen een testmatrix van $3^3 = 27$ unieke \gls{rag}-configuraties vormen (Sectie~\ref{ch:shortlist}), en de dataset en golden dataset die gebruikt worden om de componenten op te testen (Hoofdstuk~\ref{ch:data-goldenset}). Dit hoofdstuk beschrijft hoe deze configuraties technisch zijn uitgewerkt tot een dynamische proefopstelling: welke technologieën werden gekozen en waarom, hoe de individuele componenten zijn geïmplementeerd en hoe de pipeline is opgebouwd. De uitvoering van de proefopstelling en de infrastructuur waarop ze draait, komen aan bod in Hoofdstuk~\ref{ch:pipeline-testen}.

%%----------------------------------------------------------------------
\section{Haystack 2.x als orchestratieframework}
\label{sec:haystack-keuze}

De volledige pipeline is opgebouwd rond het open-source orchestratieframework Haystack 2.x \autocite{HaystackPipelines2026}. De keuze voor Haystack tegenover populaire alternatieven zoals LangChain of LlamaIndex is geen triviale beslissing: alle drie bieden abstracties voor het bouwen van \gls{rag}-systemen, maar ze verschillen in de mate van controle die de ontwikkelaar behoudt over de interne werking.

LangChain en LlamaIndex bieden high-level abstractions die in één aanroep data parsing, chunking en embedding combineren. Hoewel dit de ontwikkeltijd verlaagt, maakt het de afzonderlijke componenten onzichtbaar en moeilijk uitwisselbaar: een klassiek \emph{black-box}-probleem. Dit staat haaks op de \textbf{reproduceerbaarheid}-requirement (Tabel~\ref{tab:moscow}): wanneer een framework de interne werking verbergt, is het onmogelijk te garanderen dat identieke invoer steeds identieke verwerking oplevert, noch om te isoleren welk component verantwoordelijk is voor een bepaalde actie. Voor een benchmarkopstelling waarbij de impact van elke individuele keuze — chunker, embedder, \gls{llm} — afzonderlijk en meetbaar moet zijn, is volledige transparantie over de gegevensstroom een vereiste.

Haystack 2.x lost dit op via een architectuur gebaseerd op een gerichte acyclische graaf, zoals beschreven in de literatuurstudie (Sectie~\ref{subsubsec:componentgebaseerde-architectuur}). Elke stap in de pipeline is een onafhankelijke, geïsoleerde component met strikt gedefinieerde in- en uitvoertypen, die in Haystack \emph{sockets} worden genoemd \autocite{HaystackPipelines2026}. Componenten worden aan elkaar gekoppeld via expliciete \texttt{connect()}-aanroepen, zodat de volledige dataflow inzichtelijk is en elke component afzonderlijk vervangbaar is zonder de rest van de pipeline te raken.

\begin{figure}[H]
    \centering
    \resizebox{\textwidth}{!}{%
    \begin{tikzpicture}[
        font=\sffamily,
        >=Latex,
        node distance=8mm and 10mm,
        proc/.style={
            rectangle,
            rounded corners=2pt,
            draw=textstroke,
            fill=white,
            minimum height=10mm,
            minimum width=22mm,
            align=center,
            inner sep=4pt,
            line width=0.9pt
        },
        db/.style={
            cylinder,
            shape border rotate=90,
            aspect=0.28,
            draw=dbstroke,
            fill=dbfill,
            minimum height=14mm,
            minimum width=18mm,
            align=center,
            inner sep=3pt,
            line width=1.2pt
        },
        groupbox/.style={
            rounded corners=4pt,
            inner sep=8pt,
            line width=1pt
        },
        flow/.style={
            -{Latex[length=3mm,width=2mm]},
            draw=textstroke,
            line width=1pt
        },
        bidiflow/.style={
            {Latex[length=3mm,width=2mm]}-{Latex[length=3mm,width=2mm]},
            draw=textstroke,
            line width=1pt
        }
    ]

    % Indexing Pipeline
    \node[proc] (A) {Raw Documents};
    \node[proc, right=of A] (B) {Converter};
    \node[proc, right=of B] (C) {Cleaner};
    \node[proc, right=of C] (D) {Chunker $\times 3$};
    \node[proc, right=of D] (E) {Meta Cleaner};
    \node[proc, right=of E] (F) {Embedder $\times 3$};
    \node[proc, right=of F] (W) {Writer};

    \draw[flow] (A) -- (B);
    \draw[flow] (B) -- (C);
    \draw[flow] (C) -- (D);
    \draw[flow] (D) -- (E);
    \draw[flow] (E) -- (F);
    \draw[flow] (F) -- (W);

    % Query & Evaluation
    \node[proc, below=70mm of A] (H) {Golden Dataset};
    \node[proc, right=of H] (I) {Text Embedder};
    \node[proc, right=of I] (J) {Retriever};
    \node[proc, right=of J] (K) {Prompt Builder};
    \node[proc, right=of K] (L) {LLM $\times 3$};
    \node[proc, right=of L] (M) {Results};
    \node[proc, right=of M] (N) {RAGAS Metrics};
    \node[proc, right=of N] (O) {Leaderboard};

    \draw[flow] (H) -- (I);
    \draw[flow] (I) -- (J);
    \draw[flow] (J) -- (K);
    \draw[flow] (K) -- (L);
    \draw[flow] (L) -- (M);
    \draw[flow] (M) -- (N);
    \draw[flow] (N) -- (O);

    % Shared vector database
    \node[db] (VDB) at ($(W)!0.5!(J)$) {Qdrant};
    \draw[flow] (W) -- (VDB);
    \draw[bidiflow] (VDB) -- node[right, fill=white, inner sep=1.5pt] {vector search} (J);

    % Background group boxes
    \begin{scope}[on background layer]
        \node[groupbox, draw=idxstroke, fill=idxfill, fit=(A)(B)(C)(D)(E)(F)(W),
            label={[font=\bfseries\sffamily, text=idxstroke]north:Indexing Pipeline $\cdot$ 15 collections}] {};
        \node[groupbox, draw=qestroke, fill=qefill, fit=(H)(I)(J)(K)(L)(M)(N)(O),
            label={[font=\bfseries\sffamily, text=qestroke]north:Query \& Evaluation $\cdot$ 27 configs}] {};
    \end{scope}

    \end{tikzpicture}%
    }
    \caption{Visuele weergave van de volledige pipeline-architectuur met centrale Qdrant-opslag.}
    \label{fig:pipeline-flow}
\end{figure}

De tweede doorslaggevende reden is het protocol voor \textbf{custom componenten}. Haystack 2.x voorziet de \texttt{@component}-decorator waarmee elke Python-klasse als een native, plug-and-play-compatibele pipeline-bouwsteen kan worden geregistreerd, mits zij een \texttt{run()}-methode implementeert met gedefinieerde uitvoertypen. Must-Have-requirement voor \textbf{informatie-extractie uit complexe bronnen} (Tabel~\ref{tab:moscow}): de pipeline moet bestaan uit begrijpbare, uitbreidbare bouwstenen die Turtle Srl zelf kan onderhouden, uitbreiden en aanpassen wanneer nieuwe documentformaten of \gls{esg}-standaarden dit vereisen. Alle drie de chunkers en de documentreiniger zijn bijgevolg als native Haystack custom componenten geïmplementeerd.

%%----------------------------------------------------------------------
\section{Documentverwerking met Unstructured}
\label{sec:documentverwerking}

In een vroeg stadium van de implementatie werd \texttt{PyPDFToDocument} — een Haystack-wrapper rond de \texttt{pypdf}-bibliotheek — gebruikt als documentconverter. Dit werkte aanvankelijk correct voor eenvoudige tekstuele \gls{pdf}-bestanden, maar bleek ontoereikend voor de heterogene dataset van Turtle Srl (Hoofdstuk~\ref{ch:data-goldenset}).

\gls{esg}-rapporten en interne bedrijfsdocumenten bevatten namelijk een combinatie van elementen die \texttt{pypdf} niet correct kan verwerken: tabellen worden als aaneengesloten tekstrijen ingelezen zonder structuurinformatie, tekst in meerdere kolommen wordt lineair gelezen van boven naar onder waardoor de leesvolgorde fout is en afbeeldingen met tekstuele inhoud worden volledig genegeerd. Dit ondermijnt direct de Must-Have-requirement voor \textbf{informatie-extractie uit complexe bronnen} (Tabel~\ref{tab:moscow}), die stipuleert dat het systeem bruikbare antwoorden moet kunnen destilleren uit heterogene \gls{pdf}- en Excel-bestanden. Bovendien biedt \texttt{pypdf} geen ondersteuning voor \texttt{.docx} en \texttt{.xlsx} — twee van de vier bestandsformaten die in de dataset van Turtle Srl voorkomen (Sectie~\ref{sec:dataset-beschrijving}).

De Unstructured \gls{api} biedt hier een oplossing voor in één geïntegreerd pakket. Unstructured is een open-source documentverwerkingsbibliotheek die ontworpen is voor het omzetten van ongestructureerde documenten naar gestructureerde tekstrepresentaties, met ondersteuning voor zevenenveertig verschillende bestandsformaten \autocite{UnstructuredDocs2024}.

%%----------------------------------------------------------------------
\subsection{Configuratie van de converter}

De Unstructured-converter wordt in de indexeringspipeline geconfigureerd via de officiële Haystack-integratie \texttt{UnstructuredFileConverter}:

\begin{listing}[H]
\begin{minted}{python}
return UnstructuredFileConverter(
  api_url=unstructured_url,
  document_creation_mode="one-doc-per-page",
  unstructured_kwargs={
    "strategy": "hi_res",
    "languages": ["ita", "eng"],
    "skip_infer_table_types": [],
    "pdf_infer_table_structure": True,
  },
)
\end{minted}
\caption[Unstructured converter-configuratie]{De Unstructured-converter ondersteunt Italiaans en Engels — conform de Should-Have-requirement voor meertalige ondersteuning (Tabel~\ref{tab:moscow}) — en gebruikt de \texttt{hi\_res}-strategie voor nauwkeurige extractie van tabellen en lay-outcomplexe documenten.}
\label{lst:unstructured-config}
\end{listing}

De parameter \mintinline{python}{document_creation_mode="one-doc-per-page"} zorgt ervoor dat de converter één Haystack \texttt{Document}-object per pagina van een \gls{pdf}-document maakt \autocite{UnstructuredDocs2024}. Dit is een bewuste keuze die direct bijdraagt aan de Could-Have-requirement voor \textbf{granulaire bronverwijzing op paginaniveau} (Tabel~\ref{tab:moscow}): door pagina's als afzonderlijke documenten te behandelen, blijft het paginanummer als metadata gegarandeerd beschikbaar voor elk tekstfragment dat later de pipeline doorloopt. Er werd niet gekozen voor de \mintinline{python}{"one-doc-per-element"} document creation mode, omdat deze resulteert in een versnipperde tekst die bestaat uit kleine fragmenten zoals losse tabelcellen of titels \autocite{UnstructuredDocs2024}. Daarnaast zou de \mintinline{python}{"one-doc-per-file"} optie geen mogelijkheid bieden tot het toekennen van de paginanummers aan de metadata van iedere chunk, aangezien deze de inhoud van een volledig document omzet naar één Haystack Document \autocite{UnstructuredDocs2024}. \mintinline{python}{"one-doc-per-page"} behoudt de logische flow van de informatie binnen een pagina en fungeert bovendien als een uniform startpunt voor de vergelijking tussen de drie verschillende chunkingstrategieën. Hierdoor kan in dit onderzoek objectief gemeten worden hoe de fixed size, recursive character en semantic chunking strategieën de pagina-brede informatie verder verwerken, zonder dat de resultaten vooraf worden beïnvloed door de visuele segmentatie van de converter.

De parameters \mintinline{python}{pdf_infer_table_structure=True} en \mintinline{python}{strategy="hi_res"} activeren de Unstructured vision-gebaseerde tabelstructuurdetectie voor \glspl{pdf}, wat nauwkeurigere extractie van meerkoloms tabellen oplevert. Voor Excel-bestanden bewaart Unstructured de tabelstructuur intrinsiek via \texttt{openpyxl}, ongeacht deze parameter, wat voorkomt dat data in kolommen lineair en foutief wordt samengevoegd.

%%----------------------------------------------------------------------
\subsection{Documentreiniging en paginamarkering}
\label{subsec:documentcleaner}

Na de conversie passeert elk document een zelfgeschreven\\\texttt{DocumentCleaner}-component die het \texttt{@component}-protocol van Haystack implementeert. Deze component vervult twee rollen: tekstnormalisatie en het samenvoegen van pagina's per brondocument met expliciete pagina-markers.

\textbf{Samenvoegen van pagina's.} Unstructured levert per pagina een afzonderlijk \texttt{Document}-object op. De \texttt{DocumentCleaner} groepeert alle pagina's van hetzelfde bestand en voegt ze samen tot één enkel \texttt{Document}-object. Deze ontwerpbeslissing heeft een directe invloed op de retrieval-kwaliteit: wanneer pagina's afzonderlijk blijven, kunnen chunkgrenzen nooit een alineaovergang tussen twee pagina's omvatten. Door deze samen te voegen, kunnen de chunkers cross-pagina-context behouden en semantisch coherentere fragmenten produceren, wat een positieve invloed heeft op de \textbf{context recall}.

\textbf{Pagina-markers.} Om de paginatraceerbaarheid na het samenvoegen te bewaren, injecteert de cleaner vóór elke alinea een marker in de vorm \texttt{<<<PAGE X>>>}. Na het chunken extraheert de \texttt{ChunkMetaCleaner} de eerste marker die in een chunk voorkomt en koppelt die als \texttt{page}-metadata aan het fragment. Dit mechanisme implementeert de Could-Have-requirement voor \textbf{granulaire bronverwijzing op paginaniveau} (Tabel~\ref{tab:moscow}) zonder afhankelijk te zijn van de manier waarop de chunker de grenzen trekt.

\textbf{Tekstnormalisatie} elimineert overbodige witruimte en lege regels. De metadata-sleutels \texttt{source} (bestandsnaam) en \texttt{page} worden genormaliseerd en consistent doorgegeven aan alle volgende componenten. Voor Excel-bestanden wordt de bladnaam als pagina-identifier gebruikt.

%%----------------------------------------------------------------------
\subsection{ChunkMetaCleaner}
\label{subsec:chunkmetacleaner}

Na het chunken passeert elke chunk de \texttt{ChunkMetaCleaner}. Deze component vervult drie taken. Ten eerste extraheert hij de eerste \texttt{<<<PAGE X>>>}-marker uit de chunkinhoud en slaat de bijbehorende paginawaarde op als \texttt{page}-metadata. Ten tweede verwijdert hij alle markers zodat de embeddingmodellen schone tekst krijgen zonder ruis. Ten derde dedupliceert hij chunks op inhoud: bij korte pagina's kan overlap tot identieke chunks leiden, wat dubbele vectoren in Qdrant zou introduceren en de retrieval-kwaliteit verlaagt.

%%----------------------------------------------------------------------
\section{Custom chunking-componenten}
\label{sec:chunkers}

Zoals vastgesteld in de shortlist (Sectie~\ref{sec:shortlist-chunking}) worden drie chunkingstrategieën geëvalueerd. Elk is geïmplementeerd als een zelfstandig Haystack \texttt{@component} met een uniforme interface, waardoor de drie chunkers volledig verwisselbaar zijn. De ontwerpkeuze om de strategieën te encapsuleren in custom componenten heeft drie concrete argumenten:

Ten eerste garandeert de uniforme interface de \textbf{reproduceerbaarheid} (Tabel~\ref{tab:moscow}): de exacte configuratieparameters van elke chunker zijn bevroren in de broncode, onafhankelijk van toekomstige framework-updates. Ten tweede maakt de modulaire opbouw de \textbf{variabele configuratieondersteuning} mogelijk die de kern van dit onderzoek vormt: het Cartesisch product van alle componenten is alleen haalbaar als elke component dezelfde interface hanteert. Ten derde biedt deze structuur de transparantie en uitbreidbaarheid die Turtle Srl heeft gevraagd in het kader van de white-box-vereiste.
%%----------------------------------------------------------------------
\subsection{Fixed-size token splitting}
\label{subsec:fixed-chunker}

De eerste chunker is een wrapper rond Haystack's \texttt{RecursiveDocumentSplitter}, geconfigureerd met uitsluitend \texttt{[" ", ""]} als scheidingstekens. Dit dwingt het algoritme om uitsluitend op woordgrenzen te splitsen — nooit op alinea- of zinsgrenzen — wat resulteert in een strikt vaste tokenvenstergrootte van 1024 tokens met een overlap van 154 tokens — circa 15\% van de venstergrootte, conform de aanbeveling uit Sectie~\ref{subsubsec:chunking}.

Als deterministisch referentiepunt zonder semantisch bewustzijn is deze strategie de baseline waartegen de meerwaarde van de twee complexere strategieën wordt aangetoond, conform de benchmarkmethodologie beschreven in Sectie~\ref{sec:shortlist-chunking}.

De keuze voor \textbf{tokens} als spliteenheid — in tegenstelling tot woorden in de oorspronkelijke opzet — is een bewuste aanpassing. De maximale inputlengte van de drie embeddingmodellen is gedefinieerd in tokens. Door tokens als gemeenschappelijke maat te hanteren, worden chunks gegarandeerd binnen de modellimieten gehouden ongeacht de gemiddelde woordlengte van de ingevoerde taal (Italiaans versus Engels). Bovendien meten beide andere chunkers ook in tokens, wat de chunkgrootte als experimentele variabele isoleert van de spliteenheid.

\begin{listing}[H]
\begin{minted}{python}
@component
class FixedSizeTokenSplitter:
  def __init__(self, split_length: int = 400, split_overlap: int = 50):
    self.internal_splitter = RecursiveDocumentSplitter(
      separators=[" ", ""],
      split_length=split_length,
      split_overlap=split_overlap,
      split_unit="token",
    )

  @component.output_types(documents=List[Document])
  def run(self, documents: List[Document]):
    return self.internal_splitter.run(documents=documents)
\end{minted}
\caption[FixedSizeTokenSplitter: fixed-size chunking als deterministisch referentiepunt]{De fixed-size chunker splitst uitsluitend op woordgrenzen in een venster van 1024 tokens (overlap 154). Door \texttt{[" ", ""]} als separatoren te gebruiken binnen de \texttt{RecursiveDocumentSplitter}, is de vergelijking met de recursieve chunker volledig gecontroleerd: het enige verschil is de hiërarchie van scheidingstekens.}
\end{listing}

De Python-standaardwaarden in de constructor (\texttt{split\_length=400}, \texttt{split\_overlap=50}) dienen als terugvalwaarden; tijdens het experiment worden altijd de via HyPSTER-configuratie doorgegeven waarden van 1024 tokens en 154 tokens overlap gebruikt (Sectie~\ref{sec:hypster-config}).

%%----------------------------------------------------------------------
\subsection{Recursive character splitting}
\label{subsec:recursive-chunker}

De tweede chunker maakt gebruik van dezelfde \texttt{RecursiveDocumentSplitter} als de fixed-size chunker, maar met een volledige hiërarchie van scheidingstekens: \texttt{["\textbackslash n\textbackslash n", "\textbackslash n", ". ", " ", ""]}. Het algoritme splitst eerst op alinea-grenzen en grijpt pas recursief terug op fijnmazigere grenzen wanneer een fragment de limiet van 1024 tokens overschrijdt \autocite{LangChain2024}. Dezelfde overlap van 154 tokens als de fixed-size chunker is geconfigureerd, zodat de splitsingstrategie de enige variabele is.

Deze strategie scoort in de longlist (Sectie~\ref{sec:longlist-chunking}) het sterkst op de Must-Have-requirement voor context recall: door alinea-integriteit te bewaren, maximaliseert het de kans dat semantisch bij elkaar horende informatie — zoals de cijfers en de bijbehorende meetéénheid in een \gls{esg}-tabel — in één chunk terechtkomen en samen opgehaald worden. Dit is direct relevant voor de faithfulness-evaluatie: een \gls{llm} kan enkel correct rapporteren wat volledig in de opgehaalde context aanwezig is \autocite{Es2024}.

\begin{listing}[H]
\begin{minted}{python}
@component
class RecursiveCharacterSplitter:
  def __init__(
    self,
    separators: List[str] = ["\n\n", "\n", ". ", " ", ""],
    split_length: int = 400,
    split_overlap: int = 50
  ):
    self.internal_splitter = RecursiveDocumentSplitter(
      separators=separators,
      split_length=split_length,
      split_overlap=split_overlap,
      split_unit="token",
    )

  @component.output_types(documents=List[Document])
  def run(self, documents: List[Document]):
    return self.internal_splitter.run(documents=documents)
\end{minted}
\caption[RecursiveCharacterSplitter: hiërarchische splitsing op structurele grenzen]{Identieke tokenparameters als de fixed-size chunker (1024 tokens, overlap 154) maar met een volledige scheidingsteken-hiërarchie die alinea- en zinsgrenzen respecteert. De gecontroleerde parameteridentiteit maakt een directe prestatievergelijking van de splitsingstrategie mogelijk.}
\label{lst:recursive-chunker}
\end{listing}

Ook hier dienen de Python-standaardwaarden (\texttt{split\_length=400}, \texttt{split\_overlap=50}) als terugvalwaarden; het experiment gebruikt de door HyPSTER doorgegeven waarden van 1024 tokens en 154 tokens overlap (Sectie~\ref{sec:hypster-config}).

%%----------------------------------------------------------------------
\subsection{Semantische chunking}
\label{subsec:semantic-chunker}

De derde strategie is de semantische chunker, gebaseerd op de \emph{semantic chunking}-methodiek van~\textcite{Kamradt2023} zoals beschreven in de literatuurstudie (Sectie~\ref{subsubsec:chunking}). In tegenstelling tot structurele of tekstuele splitsing hanteert deze component betekenisgrenzen als primair criterium: een segmentatiegrens wordt geplaatst wanneer de cosine-gelijkenis tussen opeenvolgende zingroepen onder het 95e percentiel van de afstandsverdeling daalt~\autocite{Kamradt2023}.

Voor de berekening van de zin-embeddings wordt het \texttt{paraphrase-multilingual-mpnet-base-v2}-model ingezet. Dit meertalige model is bewust gekozen boven meer lichtgewicht Engelstalige alternatieven, conform de Should-Have-requirement voor \textbf{meertalige ondersteuning} (Tabel~\ref{tab:moscow}): de dataset bevat Italiaanse en Engelse documenten, waardoor een meertalig model noodzakelijk is om semantische grenzen correct te detecteren in beide talen. Per groep van twee zinnen (\texttt{sentences\_per\_group=2}) wordt een gecombineerde embedding berekend om ruis door korte zinnen te reduceren.

Een bewuste keuze in de experimentele opzet is dat de semantische chunking \textbf{niet} wordt uitgevoerd met de embeddingmodellen die deel uitmaken van de evaluatiematrix. Indien het embeddingmodel - dat als aparte variabele geëvalueerd wordt - ook de chunking zou aansturen, ontstaat er een afhankelijkheid tussen beide variabelen. Door een onafhankelijk, statisch model te gebruiken voor de segmentatie, worden de variabelen \textit{chunkingstrategie} en \textit{embedding-kwaliteit} strikt van elkaar geïsoleerd. Dit waarborgt de interne validiteit van het onderzoek: eventuele prestatieverschillen kunnen zo eenduidig worden toegeschreven aan de specifieke werking van de chunkingstrategie, onafhankelijk van het embeddingmodel dat later de uiteindelijke retrieval uitvoert.

\begin{listing}[H]
\begin{minted}{python}
@component
class SemanticEmbeddingChunker:
  def __init__(
    self,
    max_length: int,
    min_length: int,
    model_name: str = "sentence-transformers/paraphrase-multilingual-mpnet-base-v2",
    sentences_per_group: int = 2,
    percentile: float = 0.95,
    language: str = "it",
  ):
    self.max_length = max_length

    self.embedder = SentenceTransformersDocumentEmbedder(
      model=model_name,
      progress_bar=False,
    )
    self.splitter = EmbeddingBasedDocumentSplitter(
      document_embedder=self.embedder,
      sentences_per_group=sentences_per_group,
      percentile=percentile,
      min_length=min_length,
      max_length=max_length,
      language=language,
    )
\end{minted}
\caption[SemanticEmbeddingChunker: betekenisgrens-detectie]{De semantische chunker combineert een embedder en een splitter om grenzen te detecteren op basis van het 95e percentiel van de cosine-afstanden tussen zinnen.}
\label{lst:semantic-chunker}
\end{listing}

Een bijzonderheid van de implementatie zijn de \texttt{max\_length}- en \texttt{min\_length}-parameters. \texttt{max\_length} wordt voor elke indexeringsrun geïnitialiseerd op basis van het geselecteerde \gls{llm}; de modelspecifieke berekening wordt toegelicht in Sectie~\ref{sec:technische-uitvoering-chunking}. De \texttt{min\_length}-parameter stelt een minimale karakterlengte per zingroep in: zingroepen die korter zijn dan deze drempel, worden samengevoegd met de volgende groep voordat de semantische splitsing wordt toegepast. Dit voorkomt over-fragmentatie — excessief kleine chunks die de retrieval-precisie schaden. De achterliggende drempelwaarde van 128 tokens — omgezet naar modelspecifieke karakterwaarden via de formule in Sectie~\ref{sec:technische-uitvoering-chunking} en weergegeven in Tabel~\ref{tab:chunk-specs} — is gemotiveerd door \textcite{BhatEtAl2025}, die aantoont dat dit de bovengrens vormt van het bereik dat optimaal is voor feitgerichte vraag-antwoordretrieval, het primaire gebruiksscenario van de Turtle Srl-chatbot (Sectie~\ref{subsubsec:chunking}). Dit minimum is uitsluitend van toepassing op de semantische chunker; de twee deterministische chunkers — \texttt{FixedSizeTokenSplitter} en \texttt{RecursiveCharacterSplitter} — hanteren geen minimale chunklengte en beschouwen elk resultaatfragment ongeacht de grootte als geldig.

Omdat de \texttt{EmbeddingBasedDocumentSplitter} niet altijd elke chunk onder de maximale lengte kan houden — met name bij tabellen of structuurloze blokken zonder zinsgrens — is een harde terugval ingebouwd. Chunks die na de semantische splitsing toch de limiet overschrijden, worden via een \texttt{RecursiveDocumentSplitter} met woordgrenzen als vangnet verder opgedeeld. De terugval hanteert geen overlap (\texttt{split\_overlap=0}): een overlap zou de fragmenten langer maken dan de limiet en zo het doel van de terugval ondermijnen.

\begin{listing}[H]
\begin{minted}{python}
self._fallback = RecursiveDocumentSplitter(
  separators=["\n\n", "\n", ". ", " "],
  split_length=self.max_length,
  split_overlap=0,
  split_unit="char"
)

@component.output_types(documents=List[Document])
def run(self, documents: List[Document]) -> dict:
  chunks = self.splitter.run(documents=documents)["documents"]
  final: List[Document] = []
  for doc in chunks:
    if len(doc.content or "") > self.max_length:
      sub_chunks = self._fallback.run(documents=[doc])["documents"]
      for sub in sub_chunks:
        final.append(dataclasses.replace(sub, meta={**doc.meta, **sub.meta}))
    else:
      final.append(doc)
  return {"documents": final}
\end{minted}
\caption[SemanticEmbeddingChunker: fallback voor overgrote chunks]{Bij structuurloze blokken — zoals grote tabellen zonder zinsgrens — kan de \texttt{EmbeddingBasedDocumentSplitter} de maximale lengte niet altijd respecteren. De \texttt{RecursiveDocumentSplitter} als harde terugval garandeert dat elke chunk de embedding-modellimiet nooit overschrijdt.}
\label{lst:semantic-chunker-fallback}
\end{listing}

Hoewel de kernlogica voor de tekstsplitsing gebruikmaakt van de robuuste implementaties binnen het Haystack-framework, zijn deze strategieën geëncapsuleerd in specifieke custom componenten. Deze architecturale beslissing is ingegeven door drie concrete argumenten binnen de benchmark-methodologie:

Ten eerste garandeert deze abstractielaag een uniforme interface binnen de proefopstelling. Door elke strategie te wrappen in een gestandaardiseerd \texttt{@component}, kunnen de 27 unieke configuraties op een modulaire wijze door de pipeline worden gestroomlijnd zonder dat de orchestrator kennis hoeft te hebben van de onderliggende library-specifieke parameters.

Ten tweede bevordert deze encapsulatie de reproduceerbaarheid van het experiment. De custom componenten fungeren als een 'contract' dat de exacte configuratie van de Haystack-splitters bevriest. Hierdoor is de proefopstelling onafhankelijk van eventuele wijzigingen in de standaardinstellingen van toekomstige framework-updates.

Ten derde biedt deze structuur de noodzakelijke flexibiliteit voor samengestelde logica, zoals gedemonstreerd in de semantische chunker. Hierbij worden meerdere atomaire Haystack-componenten — een embedder en een splitter — gecombineerd tot één logische eenheid. Deze mate van granulariteit en controle is noodzakelijk om te voldoen aan de \emph{white-box} vereisten van dit onderzoek, waarbij elke stap in de gegevensstroom transparant en manipuleerbaar moet zijn.

%%----------------------------------------------------------------------
\section{Lokale embeddingmodellen}
\label{sec:lokale-embedding}

In de initiële opzet werden de embeddingmodellen aangestuurd via externe \gls{api}-diensten, zoals de \gls{nvidia} NIM \gls{api} en de Hugging Face Serverless Inference \gls{api}. Gedurende de implementatiefase is hiervan afgestapt om drie redenen die elk een directe link hebben met de requirements.

Ten eerste introduceerden externe aanroepen een afhankelijkheid van de beschikbaarheid van externe diensten. Storingen of rate limits konden de volledige benchmarkrun onderbreken, wat de \textbf{reproduceerbaarheid} (Tabel~\ref{tab:moscow}) ondermijnde. Ten tweede deprecieerden providers tijdens het onderzoek modellen uit hun catalogus, waardoor onvoorziene aanpassingen in de pipeline-architectuur noodzakelijk werden. Ten derde biedt lokale uitvoering volledige \textbf{data soevereiniteit}: documentinhoud verlaat de lokale machine niet voor de embedding-berekeningen, wat \gls{gdpr}-compliant dataverwerking garandeert (Tabel~\ref{tab:moscow}).

De definitieve keuze is om alle drie de embeddingmodellen volledig lokaal te draaien via de \texttt{sentence-transformers}-bibliotheek als inferentiebackend \autocite{ReimersGurevych2019}, in de Docker-container met \gls{gpu}-ondersteuning. De implementatie verschilt echter per model, wat zichtbaar is in de \texttt{\_make\_doc\_embedder()}-hulpfunctie:

\begin{listing}[H]
\begin{minted}{python}
def _make_doc_embedder(emb_cfg: dict):
  if emb_cfg["api_model"] == _BGE_M3:
    return BGEM3HybridDocumentEmbedder()
  truncate_dim = emb_cfg.get("truncate_dim")
  return SentenceTransformersDocumentEmbedder(
    model=emb_cfg["api_model"],
    prefix=emb_cfg.get("doc_prefix", ""),
    batch_size=4,
    truncate_dim=truncate_dim,
  )
\end{minted}
\caption[\_make\_doc\_embedder: modelspecifieke embedder-selectie]{BGE-M3 vereist een afwijkende implementatie omwille van zijn hybride dense+sparse uitvoer en krijgt de custom \texttt{BGEM3HybridDocumentEmbedder}. Arctic-embed en E5-instruct gaan via de standaard Haystack \texttt{SentenceTransformersDocumentEmbedder}.}
\label{lst:make-doc-embedder}
\end{listing}

%%----------------------------------------------------------------------
\subsection{Arctic-embed en E5-instruct via sentence-transformers}

Arctic-embed-l-v2.0 en multilingual-e5-large-instruct worden geïntegreerd via Haystack's \texttt{SentenceTransformersDocumentEmbedder}-component \autocite{HaystackPipelines2026}. Deze component laadt het model lokaal via de \texttt{sentence-transformers}-bibliotheek als inferentiebackend \autocite{ReimersGurevych2019}, waarna PyTorch de berekeningen uitvoert op de \gls{gpu} via \gls{cuda}-kernels. De modelgewichten worden bij de eerste run gedownload van de Hugging Face Hub naar de persistente Docker-volume cache en blijven daarna lokaal beschikbaar.

De inferentie verloopt via PyTorch: de modelspecifieke tokenizer zet de invoertekst om in tokenreeksen die als tensoren door de transformer-encoder worden geleid. Via mean pooling over alle tokenposities ontstaat één embeddingsvector van vaste lengte per invoertekst. PyTorch voert deze berekeningen uit via \gls{cuda}-kernels, wat een versnelling van factor tien tot vijftig oplevert ten opzichte van \gls{cpu}-inferentie.

Voor arctic-embed wordt via de \texttt{truncate\_dim}-parameter de Matryoshka-truncatie naar 256 dimensies geactiveerd (Sectie~\ref{subsubsec:matryoshka}). Voor E5-instruct wordt de domeinspecifieke query-instructieprefix via de \texttt{prefix}-parameter meegegeven, wat de asymmetrische query-document-modellering activeert \autocite{WangEtAl2022}.

Een bewuste ontwerpkeuze betreft de maximale inputlengte van \texttt{multilingual-e5-large-instruct}: het model verwerkt slechts 512 tokens per invoertekst \autocite{WangEtAl2024b}, terwijl chunks een maximale grootte hebben van 1.024 tokens (Sectie~\ref{sec:chunkers}). De \texttt{SentenceTransformersDocumentEmbedder} kapt invoer die deze limiet overschrijdt automatisch af. Er is bewust \textbf{niet} voor gekozen de chunkgrootte specifiek voor \texttt{E5-large} te halveren: het verkleinen van de chunkgrootte voor één embedder zou de chunkgrootte als gecontroleerde experimentele variabele doorbreken en een groter methodologisch confound introduceren dan de truncatie zelf. De benchmarkmethodologie vereist dat elke embedder evalueert onder identieke chunk-condities, zodat prestatieverschillen eenduidig aan het embeddingmodel worden toegeschreven en niet aan een ongelijke datavoorbereiding.

%%----------------------------------------------------------------------
\subsection{BGE-M3 via FlagEmbedding}
\label{subsec:bge-m3-embedder}

BGE-M3 vereist een afwijkende implementatie ten opzichte van arctic-embed en E5-instruct. Zoals beschreven in de short list (Sectie~\ref{sec:shortlist-embedding}), produceert BGE-M3 naast dense vectoren ook sparse lexicale gewichten voor hybride zoeken \autocite{Chen2024}. Haystack's standaard \texttt{SentenceTransformersDocumentEmbedder} ondersteunt enkel dense vectoren, waardoor de hybride zoekfunctionaliteit verloren zou gaan bij gebruik van die component.

Om dit op te lossen zijn twee custom Haystack-componenten geïmplementeerd via de \texttt{FlagEmbedding}-bibliotheek, één per fase van de pipeline. Het onderscheid tussen beide is inhoudelijk relevant: BGE-M3 hanteert bewust verschillende encoderingsstrategieën voor documenten en queries. Via \texttt{encode\_corpus()} worden documentpassages gecodeerd voor opslag; via \texttt{encode\_queries()} worden zoekvragen gecodeerd voor retrieval. Deze asymmetrie verbetert de retrieval-kwaliteit doordat het model de representatie kan optimaliseren voor de specifieke rol van elk teksttype \autocite{Chen2024}.

\textbf{Indexeringsfase: \texttt{BGEM3HybridDocumentEmbedder}.} Via één \texttt{encode\_corpus()}-aanroep produceert de component zowel een dense inbeddingsvector als een \texttt{SparseEmbedding}-object met lexicale gewichten, compatibel met Qdrant's hybride zoekfunctionaliteit \autocite{QdrantHybrid2025}:

\begin{listing}[H]
\begin{minted}{python}
output = self._model.encode_corpus(
  batch_texts,
  return_dense=True,
  return_sparse=True,
  return_colbert_vecs=False,
)
for doc, dense_vec, lex_weights in zip(
    batch_docs, output["dense_vecs"], output["lexical_weights"]
):
  result_docs.append(dataclasses.replace(
    doc,
    embedding=dense_vec.tolist(),
    sparse_embedding=SparseEmbedding(
      indices=[int(k) for k in lex_weights.keys()],
      values=[float(v) for v in lex_weights.values()],
    ),
  ))
\end{minted}
\caption[BGEM3HybridDocumentEmbedder: dense en sparse vectoren via
encode\_corpus()]{Eén \texttt{encode\_corpus()}-aanroep produceert per
document zowel de dense inbeddingsvector als de sparse lexicale gewichten.
Qdrant slaat beide op voor hybride zoekopdrachten via \gls{rrf}
(Sectie~\ref{subsubsec:hybride-zoeken}).}
\label{lst:bge-m3-doc-embedder}
\end{listing}

\textbf{Queryfase: \texttt{BGEM3HybridTextEmbedder}.} In de querypipeline produceert de analoge \texttt{BGEM3HybridTextEmbedder} via \texttt{encode\_queries()} zowel een dense als sparse queryvector. De query-instructie wordt via \texttt{query\_instruction\_for\_retrieval} aan het model meegegeven, conform de domeinspecifieke prefix uit de HyPSTER-configuratie (Sectie~\ref{sec:hypster-config}). In de querypipeline wordt voor BGE-M3 bijgevolg de \texttt{QdrantHybridRetriever} gebruikt in plaats van de standaard \texttt{QdrantEmbeddingRetriever}.

%%----------------------------------------------------------------------
\subsection{Werking van sentence-transformers met PyTorch}

De \texttt{sentence-transformers}-bibliotheek is een Python-wrapper rond de \texttt{transformers}-bibliotheek van Hugging Face, geoptimaliseerd voor het berekenen van semantische vectorrepresentaties van teksten \autocite{ReimersGurevych2019}. Bij het laden van een model — via \texttt{SentenceTransformer("model-naam")} — wordt het gewicht van het model gedownload van de Hugging Face Hub naar de lokale cache, waarna het model in het geheugen van de \gls{gpu} wordt geladen.

De inferentie wordt afgehandeld door PyTorch via een gestructureerd proces. Eerst zet de modelspecifieke tokenizer de invoertekst om in tokenreeksen, die vervolgens als tensoren door de transformer-encoder worden geleid. De uitvoer van deze encoder bestaat uit \emph{hidden state}-vectoren (verborgen toestandsvectoren) voor elke tokenpositie. Middels een pooling-operatie (default \emph{mean pooling} voor SentenceTransformers) worden deze vectoren samengevoegd tot één enkele embeddingvector van vaste lengte per invoertekst. 

PyTorch voert deze berekeningen uit op de \gls{gpu} met behulp van \gls{cuda}-kernels. Voor de modelgroottes in dit onderzoek resulteert dit in een versnelling van factor tien tot vijftig ten opzichte van \gls{cpu}-inferentie.

Binnen Haystack worden deze modellen geïntegreerd via twee specifieke componenten: de \texttt{SentenceTransformersDocumentEmbedder} voor de indexeringspipeline en de \texttt{SentenceTransformersTextEmbedder} voor de querypipeline. Beide componenten laden het model eenmalig bij initialisatie en houden de modelgewichten gedurende de volledige looptijd van de pipeline in het \gls{gpu}-geheugen om onnodige I/O-overhead te vermijden.

%%----------------------------------------------------------------------
\section{Configuratiebeheer met HyPSTER}
\label{sec:hypster-config}

Om de 27 unieke configuraties systematisch te kunnen doorlopen, wordt HyPSTER gebruikt als configuratieframework. De configuratieruimte is opgesplitst in drie geneste deelconfiguraties. De keuze voor HyPSTER sluit direct aan bij de Should-Have-requirement voor \textbf{predefined pipelines} (Tabel~\ref{tab:moscow}): door de volledige configuratieruimte declaratief te beschrijven in één Python-functie, is elke run volledig reproduceerbaar en auditeerbaar via \texttt{on\_unknown="raise"} — elke ongeldige sleutel genereert onmiddellijk een fout.

%%----------------------------------------------------------------------
\subsection{Configuratie chunkingstrategieën}

De chunking-deelconfiguratie slaat de naam van de geselecteerde chunker op samen met de bijbehorende splitparameters: \texttt{split\_length} (1024 voor de deterministische chunkers, \texttt{None} voor semantisch) en \texttt{split\_overlap} (154 respectievelijk \texttt{None}). De werkelijke chunker-instanties worden in de indexeringspipeline aangemaakt via de \texttt{\_make\_chunker()}-hulpfunctie op basis van deze waarden. De drie opties zijn \texttt{RecursiveCharacterSplitter}, \texttt{FixedSizeTokenSplitter} en \texttt{SemanticEmbeddingChunker}.

\subsection{Configuratie embeddingmodellen}

De embedding-deelconfiguratie slaat voor elk model de modelnaam, query-instructieprefix, documentprefix, vectordimensionaliteit en eventuele \gls{mrl}-truncatiedimensie op. De instructieprefixen zijn modelspecifiek en essentieel voor correcte retrieval.

\begin{table}[h]
\centering
\small
\begin{tabular}{p{0.35\linewidth} p{0.55\linewidth}}
\toprule
\textbf{Model} & \textbf{Query-instructieprefix} \\
\midrule
\texttt{intfloat/multilingual-e5-large-instruct}
  & \texttt{Instruct: Given a question about ESG reporting and sustainability, retrieve the most relevant supporting passage\textbackslash nQuery: } \\
\texttt{Snowflake/snowflake-arctic-embed-l-v2.0}
  & \texttt{Represent this sentence for searching relevant passages: } \\
\texttt{BAAI/bge-m3}
  & \texttt{Given a question about ESG reporting and sustainability, retrieve the most relevant supporting passage: } \\
\bottomrule
\end{tabular}
\caption[Instructieprefixen per embeddingmodel]{\label{tab:embedding-prefixes} De query-instructieprefix per embeddingmodel. Documentpassages krijgen geen prefix; de asymmetrie is intentioneel en modelspecifiek \autocite{WangEtAl2022, YuEtAl2024}.}
\end{table}

De vectordimensionaliteit wordt per model vastgelegd: BGE-M3 en E5-instruct produceren 1.024-dimensionele vectoren; arctic-embed wordt via de \texttt{truncate\_dim}-parameter afgekapt naar 256 dimensies conform de Matryoshka-truncatietechniek (Sectie~\ref{subsubsec:matryoshka}).

\subsection{Configuratie \glspl{llm}}

De \gls{llm}-deelconfiguratie koppelt elke modelnaam aan een backend, \gls{api}-parameters en de modelspecifieke \texttt{max\_char\_length} en \texttt{min\_char\_length}. De achtergrond van deze karakterlimieten wordt toegelicht in Sectie~\ref{sec:technische-uitvoering-chunking}.

\begin{listing}[H]
\begin{minted}{python}
configs = {
  "Gemma-3-27b-it": {
    "backend": "hf",
    "api_model": "google/gemma-3-27b-it:scaleway",
    "api_base_url": "https://router.huggingface.co/v1",
    "max_char_length": 4347,
    "min_char_length": 543,
  },
  "Llama-3.3-70B-Instruct": {
    "backend": "hf",
    "api_model": "meta-llama/Llama-3.3-70B-Instruct:ovhcloud",
    "api_base_url": "https://router.huggingface.co/v1",
    "max_char_length": 4069,
    "min_char_length": 509,
  },
  "Mistral-Small-2603": {
    "backend": "mistral",
    "api_model": "mistral-small-2603",
    "max_char_length": 4060,
    "min_char_length": 507,
  },
}
\end{minted}
\caption[HyPSTER LLM-configuratie]{De \texttt{max\_char\_length}- en \texttt{min\_char\_length}-waarden per \gls{llm} stemmen respectievelijk de maximale en minimale chunkgrootte van de semantische chunker af op de specifieke tokenizer van het model (Sectie~\ref{sec:technische-uitvoering-chunking}). De twee \texttt{backend}-waarden bepalen welke \gls{api}-client in de querypipeline wordt geïnstantieerd: \texttt{"hf"} voor de HuggingFace Inference Router, \texttt{"mistral"} voor de officiële Mistral \gls{api}.}
\label{lst:hypster-llm}
\end{listing}

%%----------------------------------------------------------------------
\section{Indexeringspipeline}
\label{sec:indexering}

De indexeringspipeline verwerkt alle brondocumenten voor alle combinaties van chunkingstrategie en embeddingmodel. De deterministische chunkers — \texttt{RecursiveCharacterSplitter} en \texttt{FixedSizeTokenSplitter} — produceren voor elk van de drie embeddingmodellen één collectie, wat $2 \times 3 = 6$ collecties oplevert.

Voor de \texttt{SemanticEmbeddingChunker} geldt echter een afwijkende structuur. Zoals beschreven in Sectie~\ref{sec:technische-uitvoering-chunking}, wordt de \texttt{max\_length}-parameter van de semantische chunker afgestemd op de specifieke tokenizer van het geselecteerde \gls{llm}. Omdat dezelfde tekst met verschillende tokenizers tot verschillende karakterlimieten leidt, zou een gedeelde collectie per embeddingmodel deze tokenizer-afhankelijkheid verbergen en de experimentele vergelijkbaarheid ondermijnen. De semantische chunker krijgt daarom één aparte collectie \emph{per combinatie van embeddingmodel én \gls{llm}}, wat $3 \times 3 = 9$ bijkomende collecties oplevert. In totaal worden $6 + 9 = 15$ Qdrant-collecties aangemaakt.

De collectienaam wordt deterministisch gegenereerd via een hulpfunctie. Voor de deterministische chunkers is de naam \texttt{<ChunkerNaam>\_<ModelId>}; voor de semantische chunker wordt de \gls{llm}-naam als suffix toegevoegd: \texttt{<ChunkerNaam>\_<ModelId>\_<LLMNaam>}.

De scheiding per embeddingmodel is altijd noodzakelijk: hoewel BGE-M3 en E5-instruct beide 1.024-dimensionele vectoren produceren en arctic-embed met \gls{mrl}-truncatie 256-dimensionele (Sectie~\ref{subsubsec:matryoshka}), zijn de vectorruimten zelfs bij gelijke dimensionaliteit fundamenteel incompatibel. Voor BGE-M3 wordt de collectie aangemaakt met \texttt{use\_sparse\_embeddings=True} om naast de dense vector ook de sparse lexicale gewichten op te slaan die het hybride zoeken mogelijk maken (Sectie~\ref{subsubsec:hybride-zoeken}).

%%----------------------------------------------------------------------
\subsection{Geheugenbeheer bij Semantische Chunking}

De \texttt{SemanticEmbeddingChunker} laadt \texttt{paraphrase-multilingual-mpnet-base-v2} op de \gls{gpu} voor de segmentatiefase. Als het primaire embeddingmodel daarna wordt geladen, zijn beide modellen tegelijkertijd in \gls{vram}, wat kan leiden tot \gls{oom}-errors tijdens de uitvoering. De oplossing bestaat uit een tweetraps-uitvoering: na de segmentatiefase wordt het interne model van de chunker via PyTorch's \texttt{model.cpu()}-methode expliciet naar het hoofdgeheugen verplaatst en wordt \texttt{torch.cuda.empty\_cache()} aangeroepen. Pas daarna wordt het primaire embeddingmodel geladen, zodat de twee neurale netwerken nooit tegelijkertijd in \gls{vram} aanwezig zijn.

%%----------------------------------------------------------------------
\subsection{Uitvoeringslogica}

De uitvoeringslogica varieert op basis van de geselecteerde chunking-methode. Voor de \texttt{SemanticEmbeddingChunker} is de pipeline opgedeeld in twee afzonderlijke fasen. Deze fysieke scheiding is noodzakelijk om het \gls{gpu}-geheugen (\gls{vram}) tussen de segmentatie- en embedding-fase vrij te maken, waardoor \gls{oom} errors worden voorkomen.

\begin{listing}[H]
\begin{minted}{python}
if chunker_name == "SemanticEmbeddingChunker":
  chunk_pipe = Pipeline()
  chunk_pipe.add_component("converter", _make_converter(unstructured_url))
  chunk_pipe.add_component("cleaner", DocumentCleaner())
  chunk_pipe.add_component("chunker", chunker)
  chunk_pipe.add_component("meta_cleaner", ChunkMetaCleaner())
  chunk_pipe.connect("converter.documents", "cleaner.documents")
  chunk_pipe.connect("cleaner.documents",   "chunker.documents")
  chunk_pipe.connect("chunker.documents",   "meta_cleaner.documents")
  chunk_result = chunk_pipe.run({"converter": {"paths": all_file_paths}})
  chunks = chunk_result["meta_cleaner"]["documents"]

  _free_semantic_chunker_gpu(chunker)
  del chunk_pipe; gc.collect()

  embed_pipe = Pipeline()
  embed_pipe.add_component("embedder", embedder)
  embed_pipe.add_component("writer",
    DocumentWriter(document_store=document_store))
  embed_pipe.connect("embedder.documents", "writer.documents")
  embed_pipe.run({"embedder": {"documents": chunks}})
else:
  indexing_pipe = Pipeline()
  indexing_pipe.add_component("converter", _make_converter(unstructured_url))
  indexing_pipe.add_component("cleaner",   DocumentCleaner())
  indexing_pipe.add_component("chunker",   chunker)
  indexing_pipe.add_component("meta_cleaner", ChunkMetaCleaner())
  indexing_pipe.add_component("embedder",  embedder)
  indexing_pipe.add_component("writer", DocumentWriter(document_store=document_store))
  indexing_pipe.connect("converter.documents",    "cleaner.documents")
  indexing_pipe.connect("cleaner.documents",      "chunker.documents")
  indexing_pipe.connect("chunker.documents",      "meta_cleaner.documents")
  indexing_pipe.connect("meta_cleaner.documents", "embedder.documents")
  indexing_pipe.connect("embedder.documents",     "writer.documents")
  indexing_pipe.run({"converter": {"paths": all_file_paths}})
\end{minted}
\caption[Indexeringslogica: tweeledige uitvoering voor semantische chunking, lineair voor de rest]{De tweeledige uitvoering vrijwaart de \gls{gpu} van gelijktijdige belasting. In beide paden is de \texttt{ChunkMetaCleaner} aanwezig die de pagina-markers verwerkt en chunks dedupliceert (Sectie~\ref{subsec:chunkmetacleaner}).}
\end{listing}

%%----------------------------------------------------------------------
\section{Querypipeline}
\label{sec:query-pipeline}

De querypipeline genereert voor elke van de 27 configuraties een antwoord op elke vraag uit de golden dataset (Hoofdstuk~\ref{ch:data-goldenset}). De tekstembedder hanteert altijd hetzelfde model als de indexeringsfase, wat garandeert dat queryvector en documentvectoren in dezelfde vectorruimte liggen.

\subsection{Opbouw: hybride vs.\ dense retrieval}

De querypipeline heeft twee uitvoeringspaden afhankelijk van het geconfigureerde embeddingmodel. Voor BGE-M3 worden zowel dense als sparse queryvectoren geproduceerd en gecombineerd via \gls{rrf} (Sectie~\ref{subsubsec:hybride-zoeken}). Voor arctic-embed en E5-instruct volstaat een enkel dense pad:

\begin{listing}[H]
\begin{minted}{python}
if use_hybrid:  # BGE-M3
  query_pipe.add_component("text_embedder",
    BGEM3HybridTextEmbedder(
        query_instruction=emb_cfg.get("query_prefix")))
  query_pipe.add_component("retriever",
    QdrantHybridRetriever(document_store=document_store))
  query_pipe.connect("text_embedder.embedding", "retriever.query_embedding")
  query_pipe.connect("text_embedder.sparse_embedding", "retriever.query_sparse_embedding")
else:  # Arctic-embed / E5-instruct
  query_pipe.add_component("text_embedder",
    SentenceTransformersTextEmbedder(
        model=emb_cfg["api_model"],
        prefix=emb_cfg.get("query_prefix", ""),
        truncate_dim=emb_cfg.get("truncate_dim")))
  query_pipe.add_component("retriever",
    QdrantEmbeddingRetriever(document_store=document_store))
  query_pipe.connect("text_embedder.embedding", "retriever.query_embedding")
\end{minted}
\caption[Querypipeline: hybride pad voor BGE-M3, dense pad voor de overige modellen]{De \texttt{QdrantHybridRetriever} combineert dense en sparse resultaten via \gls{rrf} \autocite{QdrantHybrid2025}. De \texttt{truncate\_dim}-parameter voor arctic-embed activeert de \gls{mrl}-truncatie naar 256 dimensies (Sectie~\ref{subsubsec:matryoshka}).}
\label{lst:query-pipeline}
\end{listing}

%%----------------------------------------------------------------------
\subsection{Uitvoering en opvangen van retriever-uitvoer}

De pipeline wordt uitgevoerd met de parameter \mintinline{python}{include_outputs_from={"retriever"}}. Standaard geeft Haystack enkel de einduitvoer van de pipeline terug — in dit geval het \gls{llm}-antwoord. Door de retriever expliciet op te nemen in de tussentijdse uitvoer, worden ook de opgehaalde chunks beschikbaar in het resultaat. Deze chunks zijn noodzakelijk als \texttt{contexts}-veld voor de \gls{ragas}-evaluatie: zowel faithfulness als context recall vereisen de daadwerkelijk opgehaalde passages als invoer (Hoofdstuk~\ref{ch:resultaten}).

%%----------------------------------------------------------------------
\subsection{Promptsjabloon en bronvermelding}
\label{subsec:prompt-template}

Het promptsjabloon is de technische operationalisering van de Must-Have-requirement voor \textbf{juridische traceerbaarheid \& auditeerbaarheid} (Tabel~\ref{tab:moscow}). Het instrueert de \gls{llm} om het antwoord te baseren op de aangeleverde context, de bestandsnaam en het paginanummer expliciet te vermelden en de uitvoer als geldig \gls{json} te structureren. Het \gls{json}-formaat vereenvoudigt de gestructureerde verwerking in de evaluatiefase en maakt geautomatiseerde extractie van antwoord en bronvermelding mogelijk:

\begin{listing}[H]
\begin{minted}{python}
_PROMPT_TEMPLATE = """Answer the question in Italian based on the context. Also include
the filename and page number or page name of the document containing
the retrieved chunk, on a new line after the answer to the question.
Output only in valid JSON.
Context:
{% for doc in documents %}
  File: {{ doc.meta['source'] }}, Page: {{ doc.meta.get('page', '?') }}
  Contents: {{ doc.content }}
{% endfor %}
Question: {{question}}"""
\end{minted}
\caption[Promptsjabloon voor de querypipeline]{Het sjabloon verplicht een Italiaans antwoord in \gls{json}-formaat met expliciete bronvermelding van bestandsnaam en paginanummer. De metadata-sleutels \texttt{source} en \texttt{page} zijn door de \texttt{DocumentCleaner} en \texttt{ChunkMetaCleaner} genormaliseerd en gegarandeerd aanwezig in elke chunk (Secties~\ref{subsec:documentcleaner}--\ref{subsec:chunkmetacleaner}).}
\label{lst:prompt-template}
\end{listing}

\section{Chunk-optimalisatie voor meertalige context}
\label{sec:technische-uitvoering-chunking}
%%----------------------------------------------------------------------

De achtergrond bij de keuze voor de gehanteerde chunkgroottes — het evenwicht tussen granulariteit en contextbehoud, en de aanbevolen tokenvensters voor contextueel begrip — is uiteengezet in Sectie~\ref{subsubsec:chunking}. Bij de vergelijking van de drie \glspl{llm} moet rekening worden gehouden met de heterogeniteit van de gehanteerde tokenizers. Elke \gls{llm} maakt gebruik van een specifiek vocabulair en algoritme om ruwe tekst om te zetten naar numerieke tokens \autocite{Wu2016, Kudo2018}. Omdat de \texttt{EmbeddingBasedDocumentSplitter} de maximale chunkgrootte uitdrukt in karakters — niet in tokens — is een modelspecifieke conversie noodzakelijk om een gelijke semantische dichtheid over de drie \glspl{llm} te garanderen \autocite{HaystackEbds2026}.

De berekening vertrekt van een minimale chunkgrootte van 128 tokens en een maximale chunkgrootte van 1.024 tokens, een gemiddelde Italiaanse woordlengte van 7,2 karakters \autocite{Eurotrad2022} en de model-specifieke \textit{subword fertility} $F$ — het gemiddeld aantal tokens per woord voor Italiaans \autocite{Allal2025}. De karakterlengte $C$ per \gls{llm} volgt uit:

\begin{equation}
C = T \times \left( \frac{\text{AVG}_{cpw}}{F} \right)
\end{equation}

Tabel~\ref{tab:chunk-specs} toont de resulterende waarden per model voor zowel de minimale als de maximale chunkgrootte.

\begin{table}[h]
\centering
\small
\begin{tabular}{lccc}
\toprule
\textbf{Model} & \textbf{Fertility ($F$)} & \textbf{$C_{\min}$ (128 tok.)} & \textbf{$C_{\max}$ (1024 tok.)} \\
\midrule
Llama-3.3-70B-Instruct   & 1,812 & 509 & 4.069 \\
Gemma-3-27b-it           & 1,696 & 543 & 4.347 \\
Mistral-Small-2603       & 1,816 & 507 & 4.060 \\
\bottomrule
\end{tabular}
\caption[Model-specifieke chunk-configuraties voor Italiaanse tekst]{\label{tab:chunk-specs}Model-specifieke chunk-configuraties voor Italiaanse tekstverwerking bij een minimale chunkgrootte van 128 tokens en een maximale chunkgrootte van 1.024 tokens.}
\end{table}

Deze \texttt{min\_char\_length}- en \texttt{max\_char\_length}-waarden worden opgeslagen in de HyPSTER \gls{llm}-configuratie (Sectie~\ref{sec:hypster-config}) en doorgegeven aan de \texttt{SemanticEmbeddingChunker} bij de instantiëring in de indexeringspipeline. Voor de twee deterministische chunkers — \texttt{FixedSizeTokenSplitter} en \texttt{RecursiveCharacterSplitter} — is deze conversie niet nodig, omdat zij native in tokens splitsen via de \texttt{RecursiveDocumentSplitter} \autocite{HaystackRecusriveSplitter2026}.

%%----------------------------------------------------------------------
\section{Resultaatstructuur}
\label{sec:resultaat-structuur}

Voor elke vraag in de golden dataset wordt een gedetailleerd resultaatwoordenboek samengesteld. Dit woordenboek bevat alle informatie die nodig is voor de latere evaluatie met \gls{ragas} en voor de vergelijkende statistische analyse:

\begin{itemize}
  \item \textbf{Identificatoren}: het \texttt{question\_id} en de originele vraagstelling.
  \item \textbf{Grondwaarheid}: \texttt{ground\_truth}, \texttt{expected\_source}, \texttt{source\_page} en\\\texttt{reference\_contexts} uit de golden dataset, die als referentie dienen voor de \gls{ragas}-metrics.
  \item \textbf{Configuratie-informatie}: de namen van de gebruikte chunker, embedder en \gls{llm}, afzonderlijk engecombineerd als \texttt{Configuration}-label.
  \item \textbf{Gegenereerd antwoord}: de volledige antwoordtekst zoals teruggegeven door de \gls{llm}.
  \item \textbf{Opgehaalde contexten}: de tekstinhoud van de chunks die door de retriever werden geselecteerd, als lijstvan tekenreeksen.
  \item \textbf{Operationele meetwaarden}: \texttt{latency} in seconden, het aantal \texttt{prompt\_tokens} en \texttt{completion\_tokens}.
\end{itemize}

De latency omvat de volledige end-to-end responstijd: van het indienen van de vraag tot het ontvangen van het \gls{llm}-antwoord, inclusief de embedding-\gls{api}-aanroep, de Qdrant-zoekopdracht en de \gls{llm}-\gls{api}-aanroep. Dit maakt latency een realistische maatstaf voor de operationele bruikbaarheid van elke configuratie in een productieomgeving.
%%----------------------------------------------------------------------